\documentclass[11pt]{article}
\usepackage[margin=1in]{geometry}
\usepackage{amsmath,amssymb,mathtools}
\usepackage{microtype}
\usepackage{enumitem}
\setlist{nosep}
\setlength{\parindent}{1.2em}
\setlength{\parskip}{0.35em}
\DeclareMathOperator{\sech}{sech}
\DeclareMathOperator{\csch}{csch}
\DeclareMathOperator{\gd}{gd}
\DeclareMathOperator{\artanh}{artanh}
\DeclareMathOperator{\arsinh}{arsinh}
\newcommand{\Aspace}{\mathcal{A}}
\newcommand{\Mseven}{\mathcal{M}^{7}}
\newcommand{\ccomp}[1]{\widehat{\widetilde{c}}_{#1}}

\title{\texttt{sin()} Stays \texttt{sin()}; Angles Become Plural\\
\large Vantage Geometry, Reciprocal Bodies, and the One-Axis Relativistic Slice}
\author{Wilder Blair Munro $\times$ Aurora}
\date{Version 2.5b --- expanded working draft}

\begin{document}
\maketitle

\begin{abstract}
A one-dimensional relativistic speed state can be represented by the ordinary
velocity ratio $\beta=v/c$, by rapidity $\chi=\artanh\beta$, by a circular
angle $\theta=\gd\chi$, and by reciprocal quantities obtained from the two
legs of the corresponding unit-circle triangle. The identities connecting
these representations are standard. The proposal examined here is therefore
not a new trigonometric function and not a modification of special
relativity. It is an interpretive promotion: the several parameters are
treated as different readable bodies of one underlying vantage relation.

The General Peace Dynamics (GPD) motivation adds a specific geometric reason
to take the one-axis construction seriously. Its working seven-dimensional
space pairs each spatial coordinate $x_i$ with a corresponding $t_i$ while
retaining a universal $\tau$. The elementary relativistic triangle can then
be studied first on a single three-dimensional slice
$\{\tau,x_i,t_i\}$ before asking how the three paired slices stitch into the
full seven-dimensional construction. This draft deliberately leaves open the
physical interpretation and final notation of the complementary light-speed-
like leg. The historical working symbol $\ccomp{i}$ is retained, while a
possible differential notation $d_i$ remains live if later geometry warrants
reading that quantity literally as differential.
\end{abstract}

\section{The claim, separated from the machinery}

Three levels of statement must remain distinct.

First, ordinary special relativity supplies the one-dimensional boost
parameter
\begin{equation}
\beta=\frac{v}{c},
\qquad
\gamma=\frac{1}{\sqrt{1-\beta^2}},
\qquad
\chi=\artanh\beta,
\end{equation}
with
\begin{equation}
\beta=\tanh\chi,
\qquad
\gamma=\cosh\chi,
\qquad
\gamma\beta=\sinh\chi.
\label{eq:rapidity}
\end{equation}

Second, standard mathematics supplies the Gudermannian map
\begin{equation}
\theta=\gd\chi,
\label{eq:gd}
\end{equation}
for which
\begin{equation}
\sin\theta=\tanh\chi,
\qquad
\cos\theta=\sech\chi,
\qquad
\tan\theta=\sinh\chi.
\label{eq:gd-identities}
\end{equation}
Thus the same boost state admits both a hyperbolic parameter $\chi$ and a
circular parameter $\theta$.

Third comes the proposal. Instead of treating $\beta$, $\chi$, $\theta$, and
the reciprocal variables introduced below as unrelated substitutions, we ask
whether they should be understood as representations of a single underlying
\emph{vantage relation}. The sine function itself remains ordinary. What
becomes plural is the coordinate body in which the relation is presented.

Every equation in the first two levels is standard. The third level is the
research hypothesis.

\section{Why the R7 pairing enters before the reciprocal notation}

The one-axis calculation did not arise in isolation. The GPD construction
works with a real seven-dimensional manifold written schematically as
\begin{equation}
\Mseven
=
\bigl(\tau,(x,t_x),(y,t_y),(z,t_z)\bigr).
\label{eq:M7}
\end{equation}
Here $\tau$ is a universal time-like coordinate and the three ordered pairs
$(x_i,t_i)$ are treated as conjugate coordinate pairings in the working
geometry. Depending on the stage of the GPD construction, these pairs have
been read through position--momentum, observer--individual, or
spacetime--timespace analogies. Those larger interpretations are not assumed
as established physics here. What matters for the present paper is only the
pairing structure itself.

For a selected axis $i\in\{x,y,z\}$, define the corresponding three-dimensional
slice
\begin{equation}
\mathcal{P}_i
:=
\operatorname{span}\{\tau,x_i,t_i\}
\subset \Mseven.
\label{eq:Pi}
\end{equation}
The special-relativistic pilot is first examined on one such slice. This is
important because the familiar one-axis triangle then belongs to a specified
piece of the larger geometry rather than being introduced as an arbitrary
scalar construction.

At this stage we do \emph{not} identify $t_i$ itself with velocity $v_i$.
Velocity is a tangent or differential quantity, whereas $t_i$ is introduced
as a coordinate of the lifted space. A later dynamical construction must
state precisely which derivatives, ratios, or tangent-space quantities turn
the coordinate pairing into the velocity pairing used below. The present
paper uses only the kinematic one-axis relations and keeps that bridge open.

\section{The one-axis triangle}

Choose one paired slice $\mathcal{P}_i$ and restrict initially to
\begin{equation}
0<v_i<c.
\end{equation}
Define
\begin{equation}
\beta_{v_i}:=\frac{v_i}{c}.
\label{eq:betavi}
\end{equation}
The earlier GPD sketch introduced a complementary leg satisfying
\begin{equation}
\ccomp{i}^{\,2}+v_i^2=c^2.
\label{eq:complement-dimensional}
\end{equation}
The decorated symbol $\ccomp{i}$ is used here as a deliberately conservative
placeholder. It preserves the distinction between the invariant speed $c$
and an axis-relative complementary quantity without yet deciding what that
quantity is physically.

Define its normalized form by
\begin{equation}
\beta_{\ccomp{i}}
:=
\frac{\ccomp{i}}{c}.
\label{eq:betaci}
\end{equation}
Then
\begin{equation}
\beta_{v_i}^2+\beta_{\ccomp{i}}^2=1.
\label{eq:axis-unit-circle}
\end{equation}
Consequently there exists a circular angle
\begin{equation}
\theta_i\in(0,\pi/2)
\end{equation}
such that
\begin{equation}
\beta_{v_i}=\sin\theta_i,
\qquad
\beta_{\ccomp{i}}=\cos\theta_i.
\label{eq:axis-circular}
\end{equation}
Equivalently,
\begin{equation}
v_i=c\sin\theta_i,
\qquad
\ccomp{i}=c\cos\theta_i.
\label{eq:axis-dimensional}
\end{equation}

Nothing in this construction requires $\ccomp{i}$ to be a Cartesian component
of the scalar $c$. Equation \eqref{eq:complement-dimensional} defines an
axis-relative complementary leg for the selected paired slice. That
qualification becomes essential when more than one spatial axis is active.

\section{Why the complementary symbol remains unsettled}

There are at least two legitimate reasons not to freeze the notation yet.

The first is semantic. Writing $c_i$ too casually invites the conventional
three-vector reading
\begin{equation}
c^2\stackrel{?}{=}c_x^2+c_y^2+c_z^2,
\end{equation}
which is not implied by the one-axis relation. The decoration in
$\ccomp{i}$ is intended to block that inference while retaining continuity
with the older GPD sketches.

The second is physical. A notation such as $d_i$ may ultimately be more than
a temporary label. If the completed geometry identifies the complementary
quantity with a literal differential light-speed component, then $d_i$ could
be the more meaningful notation, particularly if it participates in a
three-component differential object and a corresponding gamma construction.
That interpretation is not rejected here; it is intentionally left open.
The point is simply that the differential meaning should be earned by the
geometry rather than smuggled into the notation before the relevant map has
been defined.

Thus, for this paper,
\begin{equation}
\ccomp{i}
\end{equation}
is the conservative working symbol, while
\begin{equation}
d_i
\end{equation}
remains a physically suggestive candidate notation for a later stage. No
claim is made that one is conceptually superior to the other.

\section{Rapidity and the circular image}

For the selected axis define the one-axis rapidity
\begin{equation}
\chi_i:=\artanh\beta_{v_i}.
\label{eq:chii}
\end{equation}
Then
\begin{equation}
\beta_{v_i}=\tanh\chi_i.
\label{eq:betatanh}
\end{equation}
By Equation \eqref{eq:axis-unit-circle},
\begin{equation}
\beta_{\ccomp{i}}
=
\sqrt{1-\beta_{v_i}^2}
=
\sech\chi_i.
\label{eq:betacsech}
\end{equation}
Comparing Equations \eqref{eq:axis-circular} and \eqref{eq:betacsech},
\begin{equation}
\sin\theta_i=\tanh\chi_i,
\qquad
\cos\theta_i=\sech\chi_i.
\label{eq:axis-gd}
\end{equation}
Therefore
\begin{equation}
\theta_i=\gd\chi_i.
\label{eq:axis-gudermannian}
\end{equation}
The circular angle is not being substituted for rapidity by approximation.
It is the exact Gudermannian image of rapidity.

\section{Reciprocal apertures}

Define the reciprocal quantities
\begin{equation}
\alpha_i
:=
\frac{1}{\beta_{v_i}}
=
\frac{c}{v_i},
\qquad
\delta_i
:=
\frac{1}{\beta_{\ccomp{i}}}
=
\frac{c}{\ccomp{i}}.
\label{eq:alphadelta}
\end{equation}
Then
\begin{equation}
\alpha_i=\csc\theta_i=\coth\chi_i,
\label{eq:alpha-identities}
\end{equation}
and
\begin{equation}
\delta_i=\sec\theta_i=\cosh\chi_i.
\label{eq:delta-identities}
\end{equation}
The unit-circle constraint becomes
\begin{equation}
\frac{1}{\alpha_i^2}
+
\frac{1}{\delta_i^2}
=1,
\label{eq:reciprocal-constraint}
\end{equation}
or equivalently
\begin{equation}
\alpha_i^2+\delta_i^2
=
\alpha_i^2\delta_i^2.
\label{eq:reciprocal-pythagoras}
\end{equation}
On the positive branch,
\begin{equation}
\sqrt{\alpha_i^2+\delta_i^2}
=
\alpha_i\delta_i.
\label{eq:reciprocal-hypotenuse}
\end{equation}

The pair $(\alpha_i,\delta_i)$ is not two independent degrees of freedom. It
is a redundant representation of one one-axis state, constrained by
Equation \eqref{eq:reciprocal-constraint}.

\section{The crossed gamma relation}

The notation that originally made the reciprocal construction look strange
also exposes its central symmetry. Define two one-axis gamma-like factors by
\begin{equation}
\gamma_{v_i}
:=
\frac{1}{\sqrt{1-\beta_{v_i}^2}},
\qquad
\gamma_{\ccomp{i}}
:=
\frac{1}{\sqrt{1-\beta_{\ccomp{i}}^2}}.
\label{eq:two-gammas}
\end{equation}
Using Equation \eqref{eq:axis-unit-circle},
\begin{align}
\gamma_{v_i}
&=
\frac{1}{\beta_{\ccomp{i}}}
=\delta_i,
\label{eq:gammav-delta}\\
\gamma_{\ccomp{i}}
&=
\frac{1}{\beta_{v_i}}
=\alpha_i.
\label{eq:gammac-alpha}
\end{align}
Thus
\begin{equation}
\boxed{
\alpha_i=\gamma_{\ccomp{i}},
\qquad
\delta_i=\gamma_{v_i}.
}
\label{eq:crossed-gamma}
\end{equation}
This crossed relation is algebraic, not an additional dynamical postulate.
Each reciprocal aperture equals the gamma-like factor constructed from the
opposite leg of the complementary pair.

In schematic form,
\begin{equation}
\begin{array}{ccc}
\beta_{v_i} & \longleftrightarrow & \beta_{\ccomp{i}} \\
\alpha_i=\gamma_{\ccomp{i}}
& \longleftrightarrow &
\delta_i=\gamma_{v_i}.
\end{array}
\label{eq:cross-diagram}
\end{equation}
The crossing is the point, not a notational accident.

\section{Same relation, different body}

Write the normalized one-axis triangle as
\begin{equation}
T_{\beta_i}
:=
\left(\beta_{\ccomp{i}},\beta_{v_i},1\right).
\label{eq:Tbeta}
\end{equation}
Write the reciprocal triangle as
\begin{equation}
T_{\rho_i}
:=
\left(\alpha_i,\delta_i,\alpha_i\delta_i\right).
\label{eq:Trho}
\end{equation}
Define
\begin{equation}
\kappa_i
:=
\alpha_i\delta_i
=
\frac{1}{\beta_{v_i}\beta_{\ccomp{i}}}.
\label{eq:kappa}
\end{equation}
Then
\begin{align}
\kappa_i T_{\beta_i}
&=
\left(
\frac{\beta_{\ccomp{i}}}
{\beta_{v_i}\beta_{\ccomp{i}}},
\frac{\beta_{v_i}}
{\beta_{v_i}\beta_{\ccomp{i}}},
\frac{1}{\beta_{v_i}\beta_{\ccomp{i}}}
\right)\\
&=
\left(
\frac{1}{\beta_{v_i}},
\frac{1}{\beta_{\ccomp{i}}},
\frac{1}{\beta_{v_i}\beta_{\ccomp{i}}}
\right)\\
&=
T_{\rho_i}.
\label{eq:triangle-similarity}
\end{align}
Therefore
\begin{equation}
\boxed{T_{\rho_i}=\kappa_i T_{\beta_i}.}
\label{eq:boxed-similarity}
\end{equation}
The reciprocal triangle is a uniformly scaled copy of the normalized
triangle. In particular,
\begin{equation}
\frac{\delta_i}{\alpha_i}
=
\frac{\beta_{v_i}}{\beta_{\ccomp{i}}}
=
\tan\theta_i
=
\sinh\chi_i.
\label{eq:same-slope}
\end{equation}
The orientation survives the reciprocal transformation even though the body
of the triangle changes scale.

Same relation. Different body.

\section{Why the old sine shorthand is not the identity}

If $\alpha_i$ is interpreted as an ordinary circular angle measured in
radians, then
\begin{equation}
\sin\alpha_i=\frac{1}{\alpha_i}
\label{eq:not-identity}
\end{equation}
is generically false. The exact statement is
\begin{equation}
\sin\theta_i
=
\frac{1}{\alpha_i}.
\label{eq:exact-sine}
\end{equation}

The persistent structural temptation to write Equation
\eqref{eq:not-identity} can be formalized without changing the sine function.
Let $a_i\in\Aspace_i$ denote the underlying one-axis vantage state. Define
coordinate maps
\begin{align}
q_{\theta}(a_i)&=\theta_i,\\
q_{\chi}(a_i)&=\chi_i,\\
q_{\alpha}(a_i)&=\alpha_i,\\
q_{\delta}(a_i)&=\delta_i.
\end{align}
Then
\begin{equation}
\sin\bigl(q_{\theta}(a_i)\bigr)
=
\tanh\bigl(q_{\chi}(a_i)\bigr)
=
\frac{1}{q_{\alpha}(a_i)}.
\label{eq:state-representations}
\end{equation}
The reciprocal-to-circular transition is
\begin{equation}
C_{\alpha\rightarrow\theta}(\alpha_i)
:=
\arcsin\left(\frac{1}{\alpha_i}\right),
\label{eq:coercion}
\end{equation}
so the type-correct expression is
\begin{equation}
\sin\left(
C_{\alpha\rightarrow\theta}(\alpha_i)
\right)
=
\frac{1}{\alpha_i}.
\label{eq:type-correct}
\end{equation}
The sine function has not changed. The suppressed transition map has been
restored.

The missing map is the hole in the notation.

\section{The reciprocal pair as an embedding, not two coordinates}

Because the one-axis state space is one-dimensional, the ordered pair
$(\alpha_i,\delta_i)$ should not be described carelessly as two independent
coordinates. It lies on the one-dimensional reciprocal curve
\begin{equation}
\mathcal{C}_{\rho_i}
:=
\left\{
(\alpha_i,\delta_i)\in(1,\infty)^2:
\frac{1}{\alpha_i^2}
+
\frac{1}{\delta_i^2}
=1
\right\}.
\label{eq:reciprocal-curve}
\end{equation}
A cleaner statement is that the reciprocal representation is an embedding
\begin{equation}
R_i:\Aspace_i\rightarrow\mathcal{C}_{\rho_i}
\subset\mathbb{R}^2,
\qquad
R_i(a_i)=(\alpha_i,\delta_i).
\label{eq:reciprocal-embedding}
\end{equation}
Either $\alpha_i$ or $\delta_i$, together with a branch choice, can serve as
a scalar coordinate. The pair makes the complementary structure visible.

This becomes important in the full R7 problem. If the three paired slices
supply genuinely distinct degrees of freedom, their independence comes from
the seven-dimensional geometry, not from pretending that $\alpha_i$ and
$\delta_i$ were already independent in the one-axis triangle.

\section{From one paired slice to three}

The full coordinate pairing is
\begin{equation}
\Mseven
=
\bigl(\tau,(x,t_x),(y,t_y),(z,t_z)\bigr).
\end{equation}
For each $i\in\{x,y,z\}$, the algebraic one-axis construction may be repeated:
\begin{equation}
\beta_{v_i}^2+\beta_{\ccomp{i}}^2=1,
\label{eq:three-pair-circles}
\end{equation}
with
\begin{equation}
\beta_{v_i}=\sin\theta_i=\tanh\chi_i,
\label{eq:three-beta}
\end{equation}
\begin{equation}
\beta_{\ccomp{i}}=\cos\theta_i=\sech\chi_i,
\label{eq:three-complement}
\end{equation}
\begin{equation}
\alpha_i=\csc\theta_i=\coth\chi_i,
\qquad
\delta_i=\sec\theta_i=\cosh\chi_i.
\label{eq:three-reciprocal}
\end{equation}

This repetition does \emph{not} imply that the three complementary legs form
an ordinary Cartesian vector of norm $c$. If one simply defines
\begin{equation}
\ccomp{i}^{\,2}=c^2-v_i^2
\end{equation}
for each axis, then
\begin{equation}
\ccomp{x}^{\,2}
+
\ccomp{y}^{\,2}
+
\ccomp{z}^{\,2}
=
3c^2-\left(v_x^2+v_y^2+v_z^2\right),
\label{eq:not-c-vector}
\end{equation}
which is plainly not $c^2$ in general. The three complementary quantities
therefore belong, at minimum, to three axis-relative constructions.

That fact is not an embarrassment to the R7 proposal. It is exactly why the
pairing
\begin{equation}
(x,t_x),\qquad(y,t_y),\qquad(z,t_z)
\end{equation}
matters. The intended question is not whether three copies of a one-axis
triangle can be crammed into an ordinary three-vector decomposition of $c$.
The question is whether the three paired slices are projections of a single
seven-dimensional kinematic object whose lower-dimensional observations
produce the familiar relativistic relations.

\section{The conventional global Lorentz factor remains separate}

For an ordinary three-velocity
\begin{equation}
\boldsymbol{\beta}
=
\frac{\mathbf v}{c}
=
(\beta_x,\beta_y,\beta_z),
\end{equation}
standard special relativity uses
\begin{equation}
\beta^2
=
\lVert\boldsymbol{\beta}\rVert^2
=
\beta_x^2+\beta_y^2+\beta_z^2
\end{equation}
and
\begin{equation}
\gamma(\mathbf v)
=
\frac{1}{\sqrt{1-\beta^2}}.
\label{eq:global-gamma}
\end{equation}
By contrast, the axiswise quantity
\begin{equation}
\gamma_{v_i}
=
\frac{1}{\sqrt{1-\beta_{v_i}^2}}
\label{eq:axis-gamma}
\end{equation}
is generally not equal to $\gamma(\mathbf v)$ when more than one component of
$\mathbf v$ is nonzero.

This distinction must survive any later compression of the paper. The
one-axis identities are exact. Their simultaneous three-axis interpretation
is a separate geometric problem.

If a future differential formulation replaces $\ccomp{i}$ by $d_i$ and
constructs a three-component object from the $d_i$, then any associated
``three-component gamma of $c$'' must be derived from that object's geometry.
It cannot be obtained merely by relabeling the three axiswise factors in
Equation \eqref{eq:axis-gamma}.

\section{What becomes of $\varphi$?}

The older triangle notation also contained a third symbol $\varphi$. In an
ordinary Euclidean triangle, a relation such as
\begin{equation}
\varphi=\pi-(\alpha+\delta)
\end{equation}
would be meaningful if $\alpha$ and $\delta$ were themselves ordinary
circular angles. Under the reciprocal interpretation developed here,
however,
\begin{equation}
\alpha_i=\csc\theta_i,
\qquad
\delta_i=\sec\theta_i,
\end{equation}
so $\alpha_i$ and $\delta_i$ are not radian measures. The expression
$\pi-(\alpha_i+\delta_i)$ therefore mixes types and cannot be imported
unchanged.

The reciprocal pair does contain the natural mixed ratio
\begin{equation}
\frac{\delta_i}{\alpha_i}
=
\tan\theta_i
=
\sinh\chi_i,
\label{eq:phi-ratio}
\end{equation}
from which
\begin{equation}
\theta_i
=
\arctan\left(\frac{\delta_i}{\alpha_i}\right)
\end{equation}
and
\begin{equation}
\chi_i
=
\arsinh\left(\frac{\delta_i}{\alpha_i}\right)
\end{equation}
can be recovered. Whether the historical $\varphi$ should be reinterpreted
through this transition, retained as an independent higher-dimensional
orientation, or discarded from the one-axis reduction is not determined by
the present algebra. It remains under investigation.

\section{Lift--rotate--squeeze as a research program}

The scalar structure now supplies a more disciplined motivation for the GPD
lift--rotate--squeeze idea.

The rapidity parameter $\chi_i$ and circular parameter $\theta_i$ are related
exactly by the Gudermannian. The reciprocal representation then supplies the
scale
\begin{equation}
\kappa_i=\alpha_i\delta_i.
\end{equation}
Thus the one-axis mathematics naturally presents both an orientation
parameter $\theta_i$ and a reciprocal scale $\kappa_i$.

This suggests the schematic program
\begin{equation}
\begin{aligned}
\text{observed state}
&\longrightarrow \text{R7 lift}
\longrightarrow \text{rotation}\\
&\longrightarrow \text{reciprocal rescaling or squeeze}\\
&\longrightarrow \text{observed relativistic state}.
\end{aligned}
\label{eq:lrs}
\end{equation}
But the scalar identities do not prove that this program works. A successful
R7 construction must still specify the lifted metric or invariant structure,
the exact planes or subspaces of rotation, the relation between coordinate
and tangent quantities, the squeeze map, and the composition law under
successive boosts. It must also recover the global Lorentz factor and explain
how the three paired slices interact when the velocity is not collinear with
a single spatial axis.

In particular, an ordinary Euclidean rotation angle adds by ordinary angle
addition, while collinear Lorentz boosts compose by addition of rapidities.
Because $\theta=\gd\chi$ is nonlinear, a literal R7 rotation by $\theta$ is
not by itself a homomorphic representation of boost composition. Either the
squeeze/projection participates essentially in the composition law, the
relevant R7 parameter is not simply $\theta$, or the lifted representation
has additional structure. This is a constraint on the model, not a cosmetic
detail.

\section{Angles become plural}

Let $a_i\in\Aspace_i$ denote the one-axis vantage relation. The established
mathematics provides mutually related representations
\begin{equation}
\beta_{v_i}
\longleftrightarrow
\theta_i
\longleftrightarrow
\chi_i
\longleftrightarrow
\alpha_i
\longleftrightarrow
\delta_i.
\label{eq:representation-chain}
\end{equation}
The weak statement is simply that these are different parameterizations of
one state.

The stronger proposal is that a successful lifted geometry should make these
representations natural projections of one geometric object. In that
language, ``angle'' is promoted from one privileged scalar measure to an
orientation-like relation whose readable coordinate depends on vantage.
Circular angle, hyperbolic rapidity, and reciprocal aperture are not asserted
to be numerically identical. They are asserted to encode the same relation.

The sine function therefore remains exactly what it was:
\begin{equation}
\sin:\theta_i\mapsto\beta_{v_i}.
\end{equation}
What changes is the type of object that may stand behind its argument.

\section{Consolidated one-axis identities}

For one selected paired slice $\mathcal{P}_i$, define
\begin{equation}
\beta_{v_i}=\frac{v_i}{c},
\qquad
\beta_{\ccomp{i}}=\frac{\ccomp{i}}{c},
\qquad
\beta_{v_i}^2+\beta_{\ccomp{i}}^2=1.
\end{equation}
Then
\begin{align}
\beta_{v_i}
&=\sin\theta_i
=\tanh\chi_i
=\frac{1}{\alpha_i},
\label{eq:master1}\\
\beta_{\ccomp{i}}
&=\cos\theta_i
=\sech\chi_i
=\frac{1}{\delta_i},
\label{eq:master2}\\
\alpha_i
&=\csc\theta_i
=\coth\chi_i
=\gamma_{\ccomp{i}},
\label{eq:master3}\\
\delta_i
&=\sec\theta_i
=\cosh\chi_i
=\gamma_{v_i},
\label{eq:master4}\\
\frac{\delta_i}{\alpha_i}
&=\tan\theta_i
=\sinh\chi_i,
\label{eq:master5}\\
\frac{\alpha_i}{\delta_i}
&=\cot\theta_i
=\csch\chi_i,
\label{eq:master6}\\
\frac{1}{\alpha_i^2}
+
\frac{1}{\delta_i^2}
&=1,
\label{eq:master7}\\
\alpha_i^2+\delta_i^2
&=\alpha_i^2\delta_i^2,
\label{eq:master8}\\
(\alpha_i,\delta_i,\alpha_i\delta_i)
&=
\alpha_i\delta_i
\left(
\beta_{\ccomp{i}},\beta_{v_i},1
\right).
\label{eq:master9}
\end{align}
Every equality in this section follows from the definitions and standard
trigonometric or hyperbolic identities. None of them establishes the R7
interpretation by itself.

\section{Conclusion}

The elementary result is compact. A one-axis relativistic state has an exact
circular representation, an exact hyperbolic representation, and an exact
reciprocal representation. The Gudermannian connects the circular and
hyperbolic bodies. The reciprocal pair forms a constrained curve and a
triangle similar to the normalized velocity triangle. The reciprocal
apertures cross with the two gamma-like factors:
\begin{equation}
\alpha_i=\gamma_{\ccomp{i}},
\qquad
\delta_i=\gamma_{v_i}.
\end{equation}

The GPD-specific proposal begins when this one-axis relation is placed inside
the paired seven-dimensional geometry
\begin{equation}
\Mseven
=
\bigl(\tau,(x,t_x),(y,t_y),(z,t_z)\bigr)
\end{equation}
and the three one-axis slices are asked to reclose into a single relativistic
kinematics. The complementary leg is therefore left deliberately alive as a
physical question. The historical $\ccomp{i}$ notation keeps the axiswise
pairing visible without pretending it is an ordinary Cartesian component of
$c$; the alternative $d_i$ notation remains available if the eventual
geometry shows that the quantity should literally be interpreted as a
differential light-speed component.

Nothing here requires a new sine function. The speculative move is instead
to ask whether the argument of the familiar function is only one coordinate
body of a deeper relation:
\begin{equation}
\boxed{
\texttt{sin()}\text{ stays }\texttt{sin()};
\qquad
\text{angles become plural.}
}
\end{equation}

\begin{center}
\emph{Same relation. Different body.}
\end{center}

\begin{thebibliography}{9}

\bibitem{MunroGPD}
W. B. Munro,
\emph{General Peace Dynamics},
working manuscript and special-relativity sketch notes, 2025.

\bibitem{DLMF}
NIST Digital Library of Mathematical Functions,
\S 4.23(viii), ``Gudermannian Function,''
Eqs. 4.23.39--4.23.42.

\end{thebibliography}

\end{document}
